\section{Analysis of the subject}
\label{sec:analysis}

\subsection{Requirements}

\Cref{tab:requirements} lists what the subject asks for and where each item is realized. Every
requirement is met; the last column of the table doubles as a map of the source code.

\begin{table}[H]
  \centering\small
  \begin{tabularx}{\textwidth}{@{}l L l@{}}
    \toprule
    \S & Requirement & Realized by \\
    \midrule
    2.1 & Recursive array whose dimension is chosen at run time & \code{DynamicArrayND} \\
    2.1 & A class for cells; cells outside the grid are always off & \code{Cell}, \code{VoidCell} \\
    2.2 & Predefined neighborhoods G0, G2, G4, G8, G6, G26 and G$k$* & \code{Neighborhoods} \\
    2.2 & A neighborhood is a rule for traversing cells (\code{Iterable}) & \code{Neighborhood} \\
    2.2 & User-defined neighborhoods given as offsets & \code{NeighborhoodCatalog} \\
    2.3 & The eleven operators, freely nested; a class hierarchy & \code{rule.expression} \\
    2.4 & Display of 1D and 2D automata; 2D slice of 3D+ automata & \code{ui.swing}, \code{Slice} \\
    2.5 & Input file: dimension, size, slice, neighborhoods, rule, initial cells or
          \code{RANDOM(k)} & \code{config} \\
    2.6 & Menu of the input files of the current directory & \code{LauncherFrame} \\
    2.6 & Files for the Game of Life and the Sierpi\'nski triangle & \code{examples/} \\
    2.6.1 & Bonus: a fractal or chaotic rule of radius 2 & \code{explore} \\
    \bottomrule
  \end{tabularx}
  \caption{Requirements of the subject and the code that realizes them.}
  \label{tab:requirements}
\end{table}

\subsection{Decisions on points left open}

The subject leaves several questions open. Each was settled explicitly, and the answer is
documented where the user will look for it.

\paragraph{What does \code{COMPTER} count?} The subject says ``the number of cells in state 1''.
The rule language, however, can produce any integer (through \code{ADD}, for instance), so a
grid may hold other states. \code{COMPTER} was implemented to the letter: it counts cells whose
state is exactly 1. This makes automata with more than two states possible, such as Brian's
Brain where state 1 means ``firing''. Since a cell in state 2 is then invisible to
\code{COMPTER(G0)}, one operator was added, \code{ETAT()}, which returns the state of the cell
itself. It is the only extension to the language and it is redundant for two-state automata.

\paragraph{In which dimension is a predefined neighborhood valid?} The subject states that G2
is ``valid in 1D'', G4 and G8 ``in 2D'', G6 and G26 ``in 3D''. Using G8 on a 3D grid is therefore
an error, reported to the user with the list of valid names. The definition was generalized
rather than restricted (\cref{sec:neighborhoods}): for any dimension $d$, G$2d$ and G$(3^d-1)$
exist, which gives the six names of the subject for $d \le 3$ and sensible ones beyond.

\paragraph{Which display for one dimension?} A one-dimensional grid is a single row of pixels.
The figures of the subject juxtapose the successive states, and this is the display adopted:
time runs down the window.

\paragraph{What if the slice is missing?} The subject requires a file to be \emph{able} to give
the displayed plane of a grid of three dimensions or more. When a file does not, the first two
axes are displayed and the others are pinned at their middle; the plane can then be moved from
the window.

\paragraph{The orientation of coordinates.} Coordinates are array indices from 0. In a picture
the first free axis runs down (rows) and the second to the right (columns), the usual convention
for matrices and the one of Matlab, which the subject refers to for slices.

\paragraph{The example of section 2.3.} The worked example uses G8, which includes the cell
itself, although its text speaks of neighbors. The formula is reproduced verbatim in
\code{examples/subject-example.xml}, and the Game of Life uses G8* as its definition demands.

\subsection{Beyond the subject}

A few features were added because the architecture made them nearly free, and because they make
the program pleasant to use and to verify: a periodic boundary condition, the detection of
convergence and of cycles (both phenomena are described in section 1.2 of the subject), drawing
cells with the mouse, per-state colors, and a command-line interface that runs and pictures
automata without a display. The figures of this report were produced by that interface.
